% Options for packages loaded elsewhere
\PassOptionsToPackage{unicode}{hyperref}
\PassOptionsToPackage{hyphens}{url}
\documentclass[
  ignorenonframetext,
]{beamer}
\newif\ifbibliography
\usepackage{pgfpages}
\setbeamertemplate{caption}[numbered]
\setbeamertemplate{caption label separator}{: }
\setbeamercolor{caption name}{fg=normal text.fg}
\beamertemplatenavigationsymbolsempty
% remove section numbering
\setbeamertemplate{part page}{
  \centering
  \begin{beamercolorbox}[sep=16pt,center]{part title}
    \usebeamerfont{part title}\insertpart\par
  \end{beamercolorbox}
}
\setbeamertemplate{section page}{
  \centering
  \begin{beamercolorbox}[sep=12pt,center]{section title}
    \usebeamerfont{section title}\insertsection\par
  \end{beamercolorbox}
}
\setbeamertemplate{subsection page}{
  \centering
  \begin{beamercolorbox}[sep=8pt,center]{subsection title}
    \usebeamerfont{subsection title}\insertsubsection\par
  \end{beamercolorbox}
}
% Prevent slide breaks in the middle of a paragraph
\widowpenalties 1 10000
\raggedbottom
\AtBeginPart{
  \frame{\partpage}
}
\AtBeginSection{
  \ifbibliography
  \else
    \frame{\sectionpage}
  \fi
}
\AtBeginSubsection{
  \frame{\subsectionpage}
}
\usepackage{iftex}
\ifPDFTeX
  \usepackage[T1]{fontenc}
  \usepackage[utf8]{inputenc}
  \usepackage{textcomp} % provide euro and other symbols
\else % if luatex or xetex
  \usepackage{unicode-math} % this also loads fontspec
  \defaultfontfeatures{Scale=MatchLowercase}
  \defaultfontfeatures[\rmfamily]{Ligatures=TeX,Scale=1}
\fi
\usepackage{lmodern}
\ifPDFTeX\else
  % xetex/luatex font selection
\fi
% Use upquote if available, for straight quotes in verbatim environments
\IfFileExists{upquote.sty}{\usepackage{upquote}}{}
\IfFileExists{microtype.sty}{% use microtype if available
  \usepackage[]{microtype}
  \UseMicrotypeSet[protrusion]{basicmath} % disable protrusion for tt fonts
}{}
\makeatletter
\@ifundefined{KOMAClassName}{% if non-KOMA class
  \IfFileExists{parskip.sty}{%
    \usepackage{parskip}
  }{% else
    \setlength{\parindent}{0pt}
    \setlength{\parskip}{6pt plus 2pt minus 1pt}}
}{% if KOMA class
  \KOMAoptions{parskip=half}}
\makeatother
\usepackage{longtable,booktabs,array}
\newcounter{none} % for unnumbered tables
\usepackage{calc} % for calculating minipage widths
\usepackage{caption}
% Make caption package work with longtable
\makeatletter
\def\fnum@table{\tablename~\thetable}
\makeatother
\usepackage{graphicx}
\makeatletter
\newsavebox\pandoc@box
\newcommand*\pandocbounded[1]{% scales image to fit in text height/width
  \sbox\pandoc@box{#1}%
  \Gscale@div\@tempa{\textheight}{\dimexpr\ht\pandoc@box+\dp\pandoc@box\relax}%
  \Gscale@div\@tempb{\linewidth}{\wd\pandoc@box}%
  \ifdim\@tempb\p@<\@tempa\p@\let\@tempa\@tempb\fi% select the smaller of both
  \ifdim\@tempa\p@<\p@\scalebox{\@tempa}{\usebox\pandoc@box}%
  \else\usebox{\pandoc@box}%
  \fi%
}
% Set default figure placement to htbp
\def\fps@figure{htbp}
\makeatother
\setlength{\emergencystretch}{3em} % prevent overfull lines
\providecommand{\tightlist}{%
  \setlength{\itemsep}{0pt}\setlength{\parskip}{0pt}}
% Auto-generated by infrastructure/rendering/_pdf_latex_helpers.py::extract_math_font_preamble.
% Minimal header passed to Pandoc via -H so Beamer slides pick up the same math font
% as the combined-PDF path. Adjust the manuscript's preamble.md to change the font globally.
\usepackage{unicode-math}
\setmathfont{latinmodern-math.otf}
\usepackage{bookmark}
\IfFileExists{xurl.sty}{\usepackage{xurl}}{} % add URL line breaks if available
\urlstyle{same}
\hypersetup{
  hidelinks,
  pdfcreator={LaTeX via pandoc}}

\author{\texorpdfstring{}{}}
\date{}

\begin{document}

\section{Distributional Active Inference (DAIF): Convergence of Semantic
Topologies and Reinforcement Learning}\label{sec:daif-results}

\begin{frame}[fragile,allowframebreaks]{Distributional Active Inference
(DAIF): Convergence of Semantic Topologies and Reinforcement Learning}
\textbf{The claim of this section, precisely.} Three mathematical
objects that arose independently in three fields --- Firth's
distributional-semantics vectors {[}-@firth1957papers{]} in linguistics,
Bellemare--Dabney--Munos return distributions
{[}-@bellemare2017distributional{]} in reinforcement learning, and
Friston variational posteriors {[}-@friston2017active{]} in active
inference --- are the \emph{same} structural object seen from three
angles. Each assigns to a pair of states a \([0,1]\)-valued hom-value (a
similarity, a probability mass on a support atom, a posterior weight);
each does so because the underlying framework had to replace an
inadequate scalar summary (co-occurrence count, expected return, MAP
estimate) with a \emph{full distribution} in order to be expressive; and
each composes along chains by multiplying those hom-values, putting all
three inside the same \([0,1]\)-enriched category
(\autoref{sec:enriched-categories}). The convergence is non-trivial
because none of the three frameworks was designed with the others in
mind, yet the enriched-category axioms (identity, sub-multiplicative
composition, \(Z\)-matrix invertibility for magnitude) hold in all three
without modification. This convergence is what Akgül et al.
{[}-@akgul2026distributional{]} call \textbf{Distributional Active
Inference (DAIF)}. The repository implements the convergence
computationally; a fully categorical proof that the three instantiations
share a common enriched-category base in the strict sense (with a single
enriching monoidal base category, not merely compatible hom-value
scales) remains open and is flagged in the Limitations subsection below.

The \texttt{src/daif/core.py} implementation (tested in the project
suite) uses a \textbf{belief-weighted mean-field approximation}: rather
than maintaining a separate return distribution \(Z(s)\) for every
case-role state \(s\) (which would cost
\(\mathcal{O}(n \cdot N_{\text{atoms}})\) memory),
\texttt{push\_forward\_return(belief,\ transition\_matrix,\ ...)}
propagates a single return distribution weighted by the current
posterior \(q(s)\) over states. Formally, one step of the contraction
computes \(\mathbf{z} = R + \gamma\, T^{\top} q\) and collapses it to
the belief-weighted scalar \(\bar z = q^{\top} \mathbf{z}\); uncertainty
is then injected back via the quantile spread of the updating
distributional return. The approximation is exact in the limit of a
sharp posterior (\(q \to \delta_{s^\ast}\)) and recovers the full
per-state distributional Bellman operator in that regime. In exchange
for the \(\mathcal{O}(n)\) complexity, it buffers aleatoric and
epistemic uncertainty simultaneously and remains sample-efficient for
the linguistic-parsing proxy model used here (\textbf{implemented and
tested}).

The terminological collision between ``distributional'' in
distributional semantics and ``distributional'' in distributional RL is
not mere homonymy---it reflects a deep structural parallel. In both
domains, the core computational move is the same: \textbf{replacing
scalar summaries with full distributional representations.} In
linguistics, this means contextualizing word identities via probability
distributions (Firth's {[}-@firth1957papers{]} company-keeping
principle). In reinforcement learning, it replaces expected-value
estimates with quantile-approximated return distributions. In active
inference, it replaces point estimates of states with variational
posterior distributions. The enriched-categorical framework of
\autoref{sec:enriched-categories} provides the unifying abstraction: all
three are instances of \([0,1]\)-enriched categories where hom-values
encode distributional proximity rather than rigid identity.

This section presents the complete implementation and quantitative
results of the \texttt{src/daif/} subpackage---seven modules, 25 public
symbols, 224 automated tests---covering six major computational
contributions: push-forward returns (\autoref{sec:daif-pushforward}),
quantile TD and implicit quantile networks
(\autoref{sec:daif-quantile}), variational message passing and Bethe
free energy (\autoref{sec:daif-vmp}), policy selection and expected free
energy (\autoref{sec:daif-policy}), unifying ERP amplitude profiles
(\autoref{sec:daif-erp}), and convergence diagnostics
(\autoref{sec:daif-metrics}).
\end{frame}

\begin{frame}[fragile,allowframebreaks]{Push-Forward Returns and the
Distributional Bellman Operator}
\protect\phantomsection\label{sec:daif-pushforward}
The formal architecture of DAIF proceeds through three stages: (1)
reconstructing active inference via variational Bayesian inference on a
controlled Markov process; (2) defining a \emph{push-forward} operation
that iteratively maps latent-space trajectories to return distributions;
and (3) deriving a temporal-difference quantile-matching algorithm that
achieves active inference's sample-efficiency advantages within a
model-free computational architecture. This permits far-sighted parsing
without explicit transition modeling:

\begin{equation}
\mathbb{E}\left[\sum_{t=0}^{\infty} \gamma^t R(x_t, a_t) \mid x_0, a_0\right] = \int_{\mathcal{S}^{\mathbb{N}_+}} R \circ f \, d(\mathbf{S}_{\#} \mathbb{P}_{x_0, a_0}^{P_\pi})
\label{eq:eq-7-1}
\end{equation}

where \(\mathbf{S}_{\#}\) denotes the push-forward measure on
representation paths, \(f: \mathcal{S} \to \mathcal{X}\) is the
stochastic decoder, and \(\gamma \in (0,1)\) is the discount factor. The
\texttt{push\_forward\_return()} function in \texttt{src/daif/core.py}
computes this via an atomised categorical projection over
\(N_{\mathrm{atoms}}\) support points
\(\{z_i\}_{i=1}^{N_{\text{atoms}}}\) spanning
\([V_{\min}, V_{\max}]\)---the C51 architecture of Bellemare et al.
{[}-@bellemare2017distributional{]}:

\begin{equation}
Z(s,a) = \sum_{i=1}^{N_{\text{atoms}}} p_i(s,a) \cdot \delta_{z_i}
\label{eq:eq-7c-c51}
\end{equation}

where \(p_i(s,a)\) is the probability assigned to support point \(z_i\)
and \(\delta\) is the Dirac delta mass. In our framework, this C51
support structure is instantiated via the
\texttt{categorical\_return\_distribution()} function. The
distributional Bellman operator \(\mathcal{T}^{\pi}\), implemented
computationally via the \texttt{distributional\_bellman\_operator()}
multi-step contraction tracking method, then maps return distributions
forward:

\begin{equation}
\mathcal{T}^{\pi} Z(s,a) \stackrel{d}{=} R(s,a) + \gamma Z(S', A'), \quad S' \sim P(\cdot|s,a), \; A' \sim \pi(\cdot|S')
\label{eq:eq-7c-bellman}
\end{equation}

\textbf{Contraction (B1)}. By Bellemare et al.
{[}-@bellemare2017distributional, Theorem 1{]}, \(\mathcal{T}^{\pi}\) is
a \(\gamma\)-contraction in the supremum \(p\)-Wasserstein metric
\(\bar W_p\) on the space of return distributions ---
i.e.~\(\bar W_p(\mathcal{T}^{\pi} Z_1, \mathcal{T}^{\pi} Z_2) \le \gamma\, \bar W_p(Z_1, Z_2)\)
for any \(p \in [1, \infty)\). Hence the fixed point \(Z^\star\) is
unique and iterates \(\mathcal{T}^{\pi n} Z_0\) converge to \(Z^\star\)
at geometric rate \(\gamma^n\); this underwrites the convergence tracked
by \texttt{distributional\_bellman\_operator()} in
\texttt{src/daif/core.py}.

\textbf{Mean-field bound (B5)}. The belief-weighted step
\(\bar z = q^\top \mathbf{z}\) used by \texttt{push\_forward\_return()}
replaces the exact per-state operator with a single scalar collapse.
Assume bounded rewards \(\lVert R\rVert_\infty \le R_{\max}\) and an
entropy budget \(H[q] < \varepsilon\) nats. Then the approximation error
of this collapse is bounded by \(\gamma\,R_{\max}\,\varepsilon\) in the
induced \(W_1\) metric (units of return), since the worst-case mass
reallocation between roles is at most the entropy of \(q\), and each
reallocated unit of probability contributes at most \(\gamma R_{\max}\)
to \(W_1\). The approximation is therefore exact in the sharp-posterior
limit \(q \to \delta_{s^\star}\) (\(\varepsilon \to 0\)) and degrades at
most linearly in \(H[q]\) as the posterior diffuses.

For case-theoretic reasoning, DAIF implies a computational architecture
in which case assignment operates distributionally at every level: the
agent maintains not a single case diagram but a \emph{distribution over
case diagrams}, weighted by their posterior probability given the
observed linguistic evidence. This distributional perspective on case
assignment aligns naturally with the graded proto-role structure of
Dowty {[}-@dowty1991thematic{]}: a noun phrase distributes probability
mass across case roles, with the distribution sharpening as more
evidence accumulates.
\end{frame}

\begin{frame}[fragile,allowframebreaks]{Quantile Temporal Difference and
Implicit Quantile Networks}
\protect\phantomsection\label{sec:daif-quantile}
Rather than representing the return distribution as a fixed categorical
support (C51), the Quantile Regression DQN (QR-DQN) approach represents
it as a uniform mixture of \(N\) Dirac masses, one at each midpoint
quantile level \(\tau_j = (2j-1)/(2N)\) for \(j = 1,\dots,N\) ---
equivalently \(\tau_i = (2i+1)/(2N)\) for the 0-indexed form
\(i = 0,\dots,N-1\) used in \texttt{src/daif/quantile.py:66}. The
\texttt{quantile\_td\_update()} function implements the Huber quantile
loss:

\begin{equation}
\mathcal{L}_{\text{QR}}(\theta) = \frac{1}{N N'} \sum_{i=1}^{N} \sum_{j=1}^{N'} \rho_{\tau_i}^{\kappa}\!\left(\delta_{ij}\right)
\label{eq:eq-7c-qr}
\end{equation}

where \(\delta_{ij} = r + \gamma z_j' - z_i\) is the temporal-difference
error, and the Huber quantile loss \(\rho_{\tau}^{\kappa}\) is:

\begin{equation}
\rho_{\tau}^{\kappa}(u) = |\tau - \mathbf{1}[u < 0]| \cdot \mathcal{L}_{\kappa}(u), \quad \mathcal{L}_{\kappa}(u) = \begin{cases} \frac{1}{2}u^2 & |u| \leq \kappa \\ \kappa(|u| - \frac{\kappa}{2}) & \text{otherwise} \end{cases}
\label{eq:eq-7c-huber}
\end{equation}

The \textbf{Implicit Quantile Network} extension
(\texttt{implicit\_quantile\_network\_update()} in
\texttt{src/daif/quantile.py}) samples quantile levels
\(\tau \sim U[0,1]\) at inference time, enabling risk-distorted policy
selection via four modes:

\begin{longtable}[]{@{}
  >{\raggedright\arraybackslash}p{(\linewidth - 4\tabcolsep) * \real{0.3333}}
  >{\raggedright\arraybackslash}p{(\linewidth - 4\tabcolsep) * \real{0.3333}}
  >{\raggedright\arraybackslash}p{(\linewidth - 4\tabcolsep) * \real{0.3333}}@{}}
\caption{IQN risk distortion modes, formulas, and semantic roles in
case-assignment parsing.}\label{tbl:iqn-modes}\tabularnewline
\toprule\noalign{}
\begin{minipage}[b]{\linewidth}\raggedright
Mode
\end{minipage} & \begin{minipage}[b]{\linewidth}\raggedright
Distortion \(\psi_{\mathrm{IQN}}(\tau)\)
\end{minipage} & \begin{minipage}[b]{\linewidth}\raggedright
Semantic Role
\end{minipage} \\
\midrule\noalign{}
\endfirsthead
\toprule\noalign{}
\begin{minipage}[b]{\linewidth}\raggedright
Mode
\end{minipage} & \begin{minipage}[b]{\linewidth}\raggedright
Distortion \(\psi_{\mathrm{IQN}}(\tau)\)
\end{minipage} & \begin{minipage}[b]{\linewidth}\raggedright
Semantic Role
\end{minipage} \\
\midrule\noalign{}
\endhead
\textbf{neutral} & \(\psi_{\mathrm{IQN}}(\tau) = \tau{}\) & Standard
expected-value maximisation \\
\textbf{optimistic} &
\(\psi_{\mathrm{IQN}}(\tau) = \tau^{1/\eta_{\mathrm{IQN}}}\) & Prefers
high-return tails; suits exploratory parsers \\
\textbf{pessimistic} &
\(\psi_{\mathrm{IQN}}(\tau) = 1 - (1-\tau)^{1/\eta_{\mathrm{IQN}}}\) &
Over-weights low-return tails; conservative case disambiguation \\
\textbf{CVaR} &
\(\psi_{\mathrm{IQN}}(\tau) = \tau \cdot \alpha_{\mathrm{CVaR}}\) &
Linear tail compression (implementation default
\(\alpha_{\mathrm{CVaR}} = 0.25\)); risk-averse comprehension \\
\bottomrule\noalign{}
\end{longtable}

The four modes are implemented exactly in
\texttt{implicit\_quantile\_network\_update()}
(\texttt{src/daif/quantile.py}) with fixed
\(\eta_{\mathrm{IQN}} = 0.71\). \textbf{Convention note.} With
\(\eta = 0.71\) we have \(1/\eta \approx 1.408 > 1\), so
\(\tau^{1/\eta} < \tau\) and \(1-(1-\tau)^{1/\eta} > \tau\) on
\((0,1)\). In our implementation the distorted level \(\tau'\) is
multiplied directly into the asymmetric Huber weight (i.e.~the loss has
its positive-error weight scaled by \(\tau'\) and its negative-error
weight by \(1-\tau'\)). Under that \emph{weight-level} convention, the
``optimistic'' mode shrinks positive-error updates (making the agent
slower to revise upward on good news --- preference for the status quo
upper tail) and the ``pessimistic'' mode inflates positive-error updates
(making the agent track the lower tail more aggressively). Readers
cross-referencing Dabney et al.'s {[}-@dabney2018distributional{]}
sampling-level distortion (where the distortion is applied to the
sampling density of \(\tau\)) should note that our mode names follow the
\emph{weight-level} semantic rather than the sampling-level one; the
underlying mathematical formulas match across the two conventions but
their qualitative labels are mirror-images.

\textbf{Consistency (B2)}. Under i.i.d. TD samples, the empirical
minimiser of the quantile Huber loss \(\rho_\tau^\kappa\) converges to
the true \(\tau\)-quantile at rate \(O(N^{-1/2})\); see Dabney et al.
{[}-@dabney2018distributional, Theorem 2{]}. The Huber threshold
\(\kappa\) trades off robustness to outliers against bias near
\(\delta=0\), and our default \(\kappa=1\) recovers Dabney et al.'s
standard setting.

The \texttt{wasserstein\_return\_distance()} function computes
discrete-quantile approximations of both \(W_1\) (absolute area between
CDFs) and \(W_2\) (root-mean-square area) distances between return
distributions. \textbf{Approximation error (B4)}. For midpoint quantiles
\(\tau_i = (2i-1)/(2N)\) and a Lipschitz quantile function, the
estimator is O(\(N^{-2}\)) consistent with the continuous integral
\(W_p = (\int_0^1 |F_a^{-1}(\tau) - F_b^{-1}(\tau)|^p\,d\tau)^{1/p}\);
this is the canonical discretisation used in quantile-regression RL
(Dabney et al. {[}-@dabney2018distributional{]}) and is documented
explicitly in the docstring of \texttt{wasserstein\_return\_distance()}.
\end{frame}

\begin{frame}[fragile,allowframebreaks]{Variational Message Passing and
Bethe Free Energy}
\protect\phantomsection\label{sec:daif-vmp}
The orchestrator of the DAIF inference cycle is the
\texttt{distributional\_case\_assignment()} function in
\texttt{src/daif/inference.py}. It wraps the push-forward mapping and
Bayesian update routines to return a \texttt{DAIFResult} object
encapsulating the free-energy trajectory. During this loop, the
\texttt{variational\_message\_passing()} sub-function implements
iterative categorical belief refinement over the case-role posterior
\(q(\mathbf{c} \mid \mathbf{o})\). Starting from a uniform prior over
\(K\) case roles, each sweep applies the precision-weighted exponential
update:

\begin{equation}
q^{(t+1)}(c_k) \propto q^{(t)}(c_k) \cdot \exp\!\bigl(w_k \cdot o_k\bigr)
\label{eq:eq-7c-vmp}
\end{equation}

where \(w_k \in [0,1]\) is the \textbf{enriched morphism weight} for
role \(k\) read off from the case category of
\autoref{sec:enriched-categories} (acting as a precision on the update)
and \(o_k = \log p(o \mid c = k)\) is the log-likelihood of the
observation under role \(k\). After each update the distribution is
renormalised via softmax. The algorithm returns posterior probabilities
and posterior precisions
\(\Lambda_{\text{post}} = \Lambda_{\text{prior}} + \Lambda_{\text{lik}}\).
Convergence is declared when
\(\|q^{(t+1)} - q^{(t)}\|_1 < \epsilon = 10^{-6}\).

\textbf{Convergence (B3)}. \emph{Proposition.} For the
single-observation categorical factor graph used in case assignment, the
softmax update (\autoref{eq:eq-7c-vmp}) is a strict KL-contraction and
possesses a unique fixed point \(q^\star\); the \(L^1\) threshold
\(\varepsilon\) is reached in \(O(\log(1/\varepsilon))\) sweeps.
\emph{Sketch.} The unnormalised multiplicative update followed by
normalisation is the gradient step of a strictly convex problem
(minimising the Bethe free energy on a tree-structured factor graph);
Yedidia, Freeman and Weiss {[}-@yedidia2005constructing, §III--IV{]}
show that in the tree (cycle-free) case, belief propagation exactly
minimises the Bethe FE, so the iteration has a unique global minimiser.
The contraction rate is bounded by the ratio of likelihood precision to
prior precision. For factor graphs with loops our implementation
inherits the weaker ``approximate fixed point'' guarantee of loopy BP;
the current linguistic application uses a single observation factor per
word, which is tree-structured.

The \textbf{Bethe free energy} provides a tractable lower-bound
approximation to the variational free energy. In the mean-field
specialisation---where each observation constitutes an independent
factor with uniform variable degrees---the Bethe FE reduces to:

\begin{equation}
F_{\text{Bethe}}[\mathbf{q}] = \underbrace{-\sum_k q(c_k) \log p(c_k)}_{\text{prior fit}} + \underbrace{\sum_k q(c_k) \log q(c_k)}_{\text{belief entropy}} - \underbrace{\sum_t \sum_k q(c_k) \log p(o_t \mid c_k)}_{\text{likelihood}}
\label{eq:eq-7c-bethe}
\end{equation}

The \texttt{bethe\_free\_energy()} function in
\texttt{src/daif/inference.py} implements the full factor-graph Bethe FE
(Yedidia et al.~2001):
\(F_{\text{Bethe}} = \sum_\alpha \mathrm{KL}(b_\alpha \| f_\alpha) - \sum_i (d_i - 1) H(b_i)\),
where \(b_\alpha\) are factor beliefs, \(f_\alpha\) are factor
potentials, and \(d_i\) is the degree of variable \(i\) in the factor
graph. The equation above is the tractable mean-field limit presented
for clarity. The \textbf{expected information gain}
(\texttt{expected\_information\_gain()}) measures the mutual information
between observations and case-role assignments---the expected KL
divergence between posterior and prior, weighted by the marginal
likelihood of each candidate observation:

\begin{equation}
\text{EIG}(o^*) = \sum_{o^*} p(o^*) \, D_{\mathrm{KL}}\bigl(q(\mathbf{c} \mid o^*) \;\|\; p(\mathbf{c})\bigr)
\label{eq:eq-7c-eig}
\end{equation}

This mutual-information formulation properly accounts for the
probability of each observation, providing a principled measure of how
much a candidate word would reduce uncertainty about the current
case-role assignment on average.

\autoref{fig:daif-free-energy} visualises the Bethe free energy
landscape over six sequential word arrivals in the sentence \emph{``Der
Hund jagt die Katze schnell''}: variational free energy decreases with
each belief update cycle, and the KL decomposition shows the balance
between model complexity (\(D_{\mathrm{KL}}(q\|p)\)) and data fit
(\(-\mathbb{E}_q[\log p(o|s)]\)).

\begin{figure}
\centering
\pandocbounded{\includegraphics[keepaspectratio,alt={Variational free energy decreases during distributional case assignment. Left: measured F\^{}\{(t)\} values read directly from DAIFResult.fe\_trajectory (blue markers) with a smoothed reference envelope (grey dashed) and vertical dashed lines marking word arrivals. Right: the real decomposition F\^{}\{(t)\} = D\_\{\textbackslash mathrm\{KL\}\}(q\^{}\{(t)\}\_\{\textbackslash text\{posterior\}\}\textbackslash\textbar q\^{}\{(t)\}\_\{\textbackslash text\{pushed\}\}) - \textbackslash mathbb\{E\}\_\{q\^{}\{(t)\}\}{[}\textbackslash log p(o\textbar s){]}, plotted from DAIFResult.diagnostics{[}"kl\_trajectory"{]} and diagnostics{[}"loglik\_trajectory"{]} (not a schematic). Generated programmatically by src.visualization.daif\_plots.plot\_free\_energy\_convergence() from make\_free\_energy\_convergence\_data() in src/cognitive/figure\_data.py.}]{output/figures/daif_free_energy_convergence.png}}
\caption{Variational free energy decreases during distributional case
assignment. \textbf{Left}: measured \(F^{(t)}\) values read directly
from \protect\texttt{DAIFResult.fe\_trajectory} (blue markers) with a
smoothed reference envelope (grey dashed) and vertical dashed lines
marking word arrivals. \textbf{Right}: the \emph{real} decomposition
\(F^{(t)} = D_{\mathrm{KL}}(q^{(t)}_{\text{posterior}}\|q^{(t)}_{\text{pushed}}) - \mathbb{E}_{q^{(t)}}[\log p(o|s)]\),
plotted from
\protect\texttt{DAIFResult.diagnostics{[}"kl\_trajectory"{]}} and
\protect\texttt{diagnostics{[}"loglik\_trajectory"{]}} (not a
schematic). Generated programmatically by
\protect\texttt{src.visualization.daif\_plots.plot\_free\_energy\_convergence()}
from \protect\texttt{make\_free\_energy\_convergence\_data()} in
\protect\texttt{src/cognitive/figure\_data.py}.}\label{fig:daif-free-energy}
\end{figure}

\autoref{fig:daif-belief-trajectory} illustrates the full belief
trajectory: starting from uniform prior over NOM, ACC, DAT, INS, the
posterior sharpens monotonically as each morphologically marked word
supplies evidence. The determiner \emph{Der} signals nominative; the
transitive verb \emph{jagt} activates a valency frame expecting an
accusative object; the accusative article \emph{die} confirms NOM=Hund,
ACC=Katze. Entropy \(H[q]\) (centre panel) drops steeply at the second
word---the most informative item in this parse.

\begin{figure}
\centering
\pandocbounded{\includegraphics[keepaspectratio,alt={Case-role posterior sharpens from uniform prior as German morphology supplies evidence. DAIF belief trajectory during sequential disambiguation of ``Der Hund jagt die Katze schnell.'' Top: stacked area showing P(NOM), P(ACC), P(DAT), P(INS) evolution over six words with German morphological glosses. Middle: entropy H{[}q{]} with annotated steepest drop marking the most informative word. Bottom: uncertainty fan around the dominant-role probability, constructed as a simple proxy \textbackslash Delta = 1 - \textbackslash max\_k P(c\_k) scaled by fixed percentile multipliers; this is a visual surrogate, not a 51-quantile decomposition of the push-forward return distribution. Generated programmatically from src.visualization.daif\_plots.plot\_belief\_trajectory().}]{output/figures/daif_belief_trajectory.png}}
\caption{Case-role posterior sharpens from uniform prior as German
morphology supplies evidence. DAIF belief trajectory during sequential
disambiguation of \emph{``Der Hund jagt die Katze schnell.''}
\textbf{Top}: stacked area showing P(NOM), P(ACC), P(DAT), P(INS)
evolution over six words with German morphological glosses.
\textbf{Middle}: entropy \(H[q]\) with annotated steepest drop marking
the most informative word. \textbf{Bottom}: uncertainty fan around the
dominant-role probability, constructed as a simple proxy
\(\Delta = 1 - \max_k P(c_k)\) scaled by fixed percentile multipliers;
this is a visual surrogate, \emph{not} a 51-quantile decomposition of
the push-forward return distribution. Generated programmatically from
\protect\texttt{src.visualization.daif\_plots.plot\_belief\_trajectory()}.}\label{fig:daif-belief-trajectory}
\end{figure}
\end{frame}

\begin{frame}[fragile,allowframebreaks]{Policy Selection and Expected
Free Energy}
\protect\phantomsection\label{sec:daif-policy}
Active inference selects actions (here: next-word predictions or
syntactic commitments) by minimising \emph{expected free energy}
\(G{(\pi)}\). \texttt{G\_policy()} in \texttt{src/daif/policy.py}
implements the four-term decomposition used throughout this paper:

\begin{equation}
G(\pi) = \underbrace{-\mathbb{E}_{q(s)}[\log p(o\mid s,\pi)]}_{\text{ambiguity}\;(\mathcal{A})}
\;-\;\underbrace{\mathbb{E}_{q(s)}\bigl[\,H[p(s\mid o)]\,\bigr]}_{\text{epistemic value}\;(\mathcal{E})}
\;-\;\gamma\,\underbrace{\mathbb{E}_{q(s)}[v(s,\pi)]}_{\text{pragmatic value}\;(\mathcal{P})}
\;+\;\beta_{\mathrm{risk}}\,\underbrace{\mathrm{Var}_{Z}[R(\pi)]}_{\text{risk}\;(\mathcal{R})}
\label{eq:eq-7c-g}
\end{equation}

Each term is signed so that \emph{minimising} \(G(\pi)\) simultaneously
minimises ambiguity, maximises expected information gain, maximises
expected pragmatic utility, and penalises high-variance return
distributions. The weights are \(\gamma>0\) (pragmatic gain) and
\(\beta_{\mathrm{risk}}\ge 0\) (risk sensitivity, using the
distributional return variance produced by
\texttt{push\_forward\_return()}).

\textbf{Collapse identity (derivation).} Setting
\(v(s,\pi) = \log p(o_{\text{goal}} \mid s,\pi)\) and
\(\beta_{\mathrm{risk}} = 0\) reduces Eq. \ref{eq:eq-7c-g} to
\begin{align*}
G(\pi) \;&=\; -\mathbb{E}_{q(s)}[\log p(o \mid s,\pi)] \;-\; \mathbb{E}_{q(s)}[H[p(s \mid o)]] \;-\; \gamma\,\mathbb{E}_{q(s)}[\log p(o_{\text{goal}} \mid s,\pi)]\\
\;&=\; \underbrace{-\mathbb{E}_{q(s)}[\log p(o \mid s,\pi)]}_{\text{expected surprise}} \;+\; \underbrace{D_{\mathrm{KL}}\bigl(q(s\mid\pi)\,\|\,p(s)\bigr)}_{\text{epistemic value}} \;-\; \gamma\,\mathbb{E}_{q(s)}[\log p(o_{\text{goal}} \mid s,\pi)],
\end{align*} where the second equality uses the standard
active-inference identity
\(-\mathbb{E}_{q(s)}[H[p(s\mid o)]] = D_{\mathrm{KL}}(q(s\mid\pi) \,\|\, p(s)) + \text{const}\)
(Friston et al. {[}-@friston2017active{]}, Eq. 5), valid when the
posterior is close to the generative prior so that the
\(H[q(s\mid\pi)]\) term absorbs into the constant offset that cancels
across policies. This is the canonical three-term form of Friston et
al., confirming that our decomposition is a conservative generalisation
rather than a departure.

Policy selection follows a Boltzmann (softmax) distribution over
negative EFE:

\begin{equation}
P(\pi) = \frac{\exp(-\alpha_{\mathrm{pol}} \cdot G(\pi))}{\sum_{\pi'} \exp(-\alpha_{\mathrm{pol}} \cdot G(\pi'))}
\label{eq:eq-7c-softmax}
\end{equation}

where \(\alpha_{\mathrm{pol}} > 0\) is the inverse temperature (policy
softmax), distinct from the CVaR tail level \(\alpha_{\mathrm{CVaR}}\)
in the IQN table above. \texttt{softmax\_policy\_selection()} implements
this across an array of candidate policies;
\texttt{distributional\_epistemic\_value()} returns the epistemic
component alone, enabling decomposition of the policy gradient.

For case-theoretic reasoning, this means the agent selects the case
assignment (NOM/ACC/DAT/etc.) that simultaneously minimises surprise
(fits the observed morphological evidence), maximises information gain
(resolves ambiguity fastest), and respects risk sensitivity (avoids
high-variance parses in pessimistic mode). This provides a principled,
Bayes-optimal account of why certain parse strategies are preferred
cross-linguistically---they minimise expected free energy under the
agent's generative model.
\end{frame}

\begin{frame}[fragile,allowframebreaks]{ERP Amplitude Profiles from
Distributional Prediction Error}
\protect\phantomsection\label{sec:daif-erp}
The DAIF prediction module (\texttt{src/daif/prediction.py}) provides
two complementary prediction error measures, each appropriate for
different levels of the distributional hierarchy:

\textbf{Scalar DPE} (\texttt{distributional\_prediction\_error()}): For
point-prediction scenarios where the expected case role is known, the
precision-weighted surprisal provides a computationally efficient scalar
measure:

\begin{equation}
\mathrm{DPE}_{\text{scalar}}(c, q) = w_f \cdot \bigl(-\log q[c_{\text{expected}}]\bigr)
\label{eq:eq-7c-dpe-scalar}
\end{equation}

\textbf{Wasserstein DPE} (\texttt{wasserstein\_prediction\_error()}):
For full distributional comparisons between predicted and observed
return distributions, the precision-weighted Wasserstein-1 distance
provides the distributional measure:

\begin{equation}
\mathrm{DPE}(o, q) = w_f \cdot W_1(Z_{\text{predicted}}, Z_{\text{observed}})
\label{eq:eq-7c-dpe}
\end{equation}

where \(w_f = \mathcal{C}(A,B) \in [0,1]\) is the enriched morphism
weight (precision) of the violated case morphism (matching
\autoref{eq:pe-precision-error} in \autoref{sec:diagrammatic-cognition})
and \(W_1\) is the Wasserstein-1 distance.

\textbf{Reader's guide to the four DPE variants.} The DAIF framework
uses four distinct but related prediction-error measures, easily
confused because they share the name ``DPE'':

\begin{enumerate}
\tightlist
\item
  \(\mathrm{DPE}_{\text{scalar}}\) (\autoref{eq:eq-7c-dpe-scalar},
  \texttt{distributional\_prediction\_error()}) --- a
  \emph{point-belief} surprisal: precision-weighted cross-entropy on the
  currently-expected role. Used when the grammatically expected role is
  known.
\item
  \(\mathrm{DPE}\) (full Wasserstein; \autoref{eq:eq-7c-dpe},
  \texttt{wasserstein\_prediction\_error()}) --- a \emph{full
  distributional} mismatch between predicted and observed return
  distributions. Used for graded / distributional comparisons.
\item
  \(\mathrm{DPE}_{\text{semantic}}\) (\autoref{eq:eq-7c-dpe-semantic}
  below, input to \texttt{n400\_from\_return\_distribution()}) --- the
  \emph{first moment} of the distributional mismatch, i.e.~the absolute
  mean-return shift; tracks the N400 heuristic component.
\item
  \(\mathrm{DPE}_{\text{structural}}\)
  (\autoref{eq:eq-7c-dpe-structural} below, input to
  \texttt{p600\_from\_precision\_update()}) --- the \emph{full
  distributional} mismatch \(W_1(Z_{\text{pred}}, Z_{\text{obs}})\);
  tracks the P600 discrepancy component and coincides with
  \(\mathrm{DPE}\) in item 2.
\end{enumerate}

Items (3) and (4) decompose (2) into a mean-shift and a
full-distribution signal; item (1) is a simpler scalar surrogate for
point-belief use cases.

\textbf{ERP derivation from the Free Energy Principle (B6).} The two ERP
formulas below are \emph{derived} from a free-energy decomposition
rather than posited empirically. A new word arriving at the
case-assignment layer induces a change in variational free energy
\(\Delta F = F_{\text{post}} - F_{\text{prior}}\). Using
\(F = D_{\mathrm{KL}}\bigl(q \,\|\, p(s)\bigr) - \mathbb{E}_q[\log p(o|s)]\)
(the variational free-energy decomposition introduced in
\autoref{sec:cognitive-integration}) and splitting \(q \to q'\) into a
posterior-mean shift and a precision-sharpening component, a first-order
expansion gives \begin{equation*}
\Delta F \;=\; \underbrace{-\bigl(\mathbb{E}_{q'}[Z_{\text{obs}}] - \mathbb{E}_{q}[Z_{\text{pred}}]\bigr)}_{\text{mean-return shift}\;\approx\;\mathrm{DPE}_{\text{semantic}}} \;+\; \underbrace{\tfrac{1}{2}\,\Delta\Lambda\,\sigma_Z^{2}}_{\text{precision update}\;\approx\;\Delta\Lambda \cdot \mathrm{DPE}_{\text{structural}}} \;+\; O(\|\Delta q\|^2),
\end{equation*} where \(\sigma_Z^2\) is the return-distribution variance
(proxied at first order by \(W_1(Z_{\text{pred}}, Z_{\text{obs}})\)).
The mean-return shift is the heuristic, expectation-dominated component
that Kuperberg and Jaeger {[}-@kuperberg2016mean{]} and Li and Futrell
{[}-@li2023decomposition{]} associate with the N400; the
precision-update term is the discrepancy, structure-update component
that Rabovsky et al. {[}-@rabovsky2018modelling{]} and Li and Futrell
{[}-@li2024shallow{]} associate with the P600. Severity gating
\(S_{\text{violation}} \in \{0, 0.5, 1.0\}\) modulates both components
multiplicatively, reflecting the attenuation of prediction-error signals
when the violation is only mild. The precision on the semantic component
is \(w_c\) (the enriched morphism weight), and the amplitude calibration
\(s\) on the P600 converts the dimensionless free-energy increment into
μV. This yields the severity-gated decomposition:

\begin{align}
\mathrm{N400}(c) &= -\,\mathrm{DPE}_{\text{semantic}} \cdot w_c \cdot S_{\text{violation}} \label{eq:eq-7c-n400} \\
\mathrm{P600}(c) &= s\,\cdot\,\Delta\Lambda \cdot \mathrm{DPE}_{\text{structural}} \cdot S_{\text{violation}} \label{eq:eq-7c-p600}
\end{align}

where \(S_{\text{violation}} \in \{0, 0.5, 1.0\}\) encodes violation
severity (congruent / mild / strong), \(w_c\) is the enriched weight of
the case morphism,
\(\Delta\Lambda = \max(0, \Lambda_{\text{post}} - \Lambda_{\text{prior}})\)
is the precision-update magnitude reflecting the structural reanalysis
cost, and \(s > 0\) is a dimensionless amplitude-calibration constant
(default \(s=1\) in \texttt{p600\_from\_precision\_update()}). The
leading minus sign in Eq. \ref{eq:eq-7c-n400} follows
electrophysiological convention: larger semantic surprise yields a more
negative deflection at the N400 latency, consistent with the sign
produced by \texttt{n400\_from\_return\_distribution()} in
\texttt{src/daif/prediction.py}. The two \(\mathrm{DPE}\) flavours
appearing here are defined formally by:

\begin{align}
\mathrm{DPE}_{\text{semantic}}\;&=\;\bigl\lvert \mathbb{E}[Z_{\text{pred}}] - \mathbb{E}[Z_{\text{obs}}] \bigr\rvert
\label{eq:eq-7c-dpe-semantic}\\
\mathrm{DPE}_{\text{structural}}\;&=\;W_{1}\bigl(Z_{\text{pred}},\,Z_{\text{obs}}\bigr)
\label{eq:eq-7c-dpe-structural}
\end{align}

i.e.~\(\mathrm{DPE}_{\text{semantic}}\) is the absolute mean-return
mismatch (the heuristic component in the Li--Futrell decomposition,
tracking the N400) and \(\mathrm{DPE}_{\text{structural}}\) is the full
distributional Wasserstein-1 mismatch (the discrepancy component,
tracking the P600). Computationally, N400 amplitudes are extracted via
\texttt{n400\_from\_return\_distribution()} (which computes mean-return
mismatch scaled by precision and severity), while P600 amplitudes use
\texttt{p600\_from\_precision\_update()} (which computes
precision-update magnitude scaled by DPE and severity). This dual
decomposition directly mirrors the empirical finding of Li and Futrell
{[}-@li2023decomposition; -@li2024shallow{]}, who show that surprisal
decomposes into a \emph{heuristic} component tracking N400 and a
\emph{discrepancy} component tracking P600---precisely the semantic
vs.~structural split captured by \(\mathrm{DPE}_{\text{semantic}}\) and
\(\mathrm{DPE}_{\text{structural}}\) above.

Our framework explicitly accommodates Rabovsky et al.'s
{[}-@rabovsky2018modelling{]} finding that the N400 reflects a
probabilistic Bayesian belief update, extending it by formally equipping
the N400 semantic surprise as a distributional prediction-error over
explicit case boundaries. (The complementary neurobiological-timing
question raised by the ROSE model is addressed in the Limitations
subsection below, after the metrics discussion and before the CEREBRUM
integration.)

Ultimately, the \texttt{erp\_amplitude\_profile()} master function
aggregates these distinct computations into a complete
\texttt{ERPProfile} dataclass containing:

\begin{itemize}
\tightlist
\item
  \textbf{N400 amplitude} (µV) and \textbf{peak latency} (ms) for each
  case role
\item
  \textbf{P600 amplitude} (µV) and \textbf{peak latency} (ms) for each
  case role
\item
  \textbf{Time-series waveforms} sampled at 1 kHz over a 1100 ms epoch
  (−200 to +900 ms)
\end{itemize}

\autoref{fig:daif-erp-predictions} demonstrates predicted ERP amplitudes
across all eight case roles under three violation conditions. A key
result: because VOC carries the lowest enriched precision weight in the
illustration (\(w = 0.10\), the smallest among all eight roles), VOC→NOM
is the most structurally inadmissible transition and elicits the largest
P600; while GEN, with moderate precision weight (\(w = 0.70\)), elicits
a pronounced N400 but attenuated P600.

\begin{figure}
\centering
\pandocbounded{\includegraphics[keepaspectratio,alt={Distributional prediction error predicts graded N400/P600 amplitudes across all eight case roles. Left: illustrative Gaussian ERP waveforms for three violation conditions --- congruent (NOM→NOM), mild (ACC→NOM), and strong (VOC→NOM) --- with fixed template amplitudes (-1.5/3.0, -4.0/3.0, -7.0/6.0 μV N400/P600) and peaks at 380 ms / 600 ms (matching DEFAULT\_N400\_PEAK\_MS / DEFAULT\_P600\_PEAK\_MS in src/daif/prediction.py), shown to depict the mechanism of -- rather than a calibrated simulation. Middle: scatter of enriched weight w versus scalar DPE () for all eight case roles, computed via distributional\_prediction\_error() in src/daif/prediction.py. Right: real DAIF-predicted magnitudes --- mean N400 magnitude via n400\_from\_return\_distribution() and mean P600 via p600\_from\_precision\_update() across the eight roles (make\_erp\_prediction\_data() in src/cognitive/figure\_data.py) --- alongside literature-typical amplitude ranges from Kutas \& Federmeier {[}-@kutas2011thirty{]} (N400 magnitude 3--5 μV, P600 5--8 μV, shown with error bars). Note that the model predictions are on a dimensionless DPE scale (units of log-probability × enriched weight), whereas the literature values are calibrated in μV; the comparison is therefore qualitative (relative ordering and graded response to precision) rather than a numerical match. A future calibration step would require fitting a per-subject μV-per-nat scaling constant to empirical ERP data. Generated programmatically from src.visualization.daif\_plots.plot\_erp\_predictions().}]{output/figures/daif_erp_predictions.png}}
\caption{Distributional prediction error predicts graded N400/P600
amplitudes across all eight case roles. \textbf{Left}: illustrative
Gaussian ERP waveforms for three violation conditions --- congruent
(NOM→NOM), mild (ACC→NOM), and strong (VOC→NOM) --- with fixed template
amplitudes (\(-1.5/3.0\), \(-4.0/3.0\), \(-7.0/6.0\) μV N400/P600) and
peaks at 380 ms / 600 ms (matching
\protect\texttt{DEFAULT\_N400\_PEAK\_MS} /
\protect\texttt{DEFAULT\_P600\_PEAK\_MS} in
\protect\texttt{src/daif/prediction.py}), shown to depict the
\emph{mechanism} of \autoref{eq:eq-7c-n400}--\autoref{eq:eq-7c-p600}
rather than a calibrated simulation. \textbf{Middle}: scatter of
enriched weight \(w\) versus scalar DPE (\autoref{eq:eq-7c-dpe-scalar})
for all eight case roles, computed via
\protect\texttt{distributional\_prediction\_error()} in
\protect\texttt{src/daif/prediction.py}. \textbf{Right}: \emph{real}
DAIF-predicted magnitudes --- mean N400 magnitude via
\protect\texttt{n400\_from\_return\_distribution()} and mean P600 via
\protect\texttt{p600\_from\_precision\_update()} across the eight roles
(\protect\texttt{make\_erp\_prediction\_data()} in
\protect\texttt{src/cognitive/figure\_data.py}) --- alongside
literature-typical amplitude ranges from Kutas \& Federmeier
{[}-@kutas2011thirty{]} (N400 magnitude 3--5 μV, P600 5--8 μV, shown
with error bars). Note that the model predictions are on a
\emph{dimensionless} DPE scale (units of log-probability × enriched
weight), whereas the literature values are calibrated in μV; the
comparison is therefore qualitative (relative ordering and graded
response to precision) rather than a numerical match. A future
calibration step would require fitting a per-subject μV-per-nat scaling
constant to empirical ERP data. Generated programmatically from
\protect\texttt{src.visualization.daif\_plots.plot\_erp\_predictions()}.}\label{fig:daif-erp-predictions}
\end{figure}
\end{frame}

\begin{frame}[fragile,allowframebreaks]{Convergence Diagnostics and
Distributional Metrics}
\protect\phantomsection\label{sec:daif-metrics}
The \texttt{src/daif/metrics.py} module provides four diagnostic tools
for verifying DAIF model behaviour:

\textbf{Convergence diagnostics} (\texttt{convergence\_diagnostics()})
assess whether a free-energy trajectory \(\{F^{(t)}\}_{t=0}^{T}\) is
well-behaved. Given the free-energy sequence produced by VMP, the
function returns a dict with the eight fields below (keys match the
Python return value exactly):

\begin{longtable}[]{@{}
  >{\raggedright\arraybackslash}p{(\linewidth - 4\tabcolsep) * \real{0.3333}}
  >{\raggedright\arraybackslash}p{(\linewidth - 4\tabcolsep) * \real{0.3333}}
  >{\raggedright\arraybackslash}p{(\linewidth - 4\tabcolsep) * \real{0.3333}}@{}}
\caption{Convergence diagnostic metrics for DAIF free-energy
trajectories. Keys match \protect\texttt{convergence\_diagnostics()} in
\protect\texttt{src/daif/metrics.py}
one-to-one.}\label{tbl:convergence-metrics}\tabularnewline
\toprule\noalign{}
\begin{minipage}[b]{\linewidth}\raggedright
Key
\end{minipage} & \begin{minipage}[b]{\linewidth}\raggedright
Formula
\end{minipage} & \begin{minipage}[b]{\linewidth}\raggedright
Interpretation
\end{minipage} \\
\midrule\noalign{}
\endfirsthead
\toprule\noalign{}
\begin{minipage}[b]{\linewidth}\raggedright
Key
\end{minipage} & \begin{minipage}[b]{\linewidth}\raggedright
Formula
\end{minipage} & \begin{minipage}[b]{\linewidth}\raggedright
Interpretation
\end{minipage} \\
\midrule\noalign{}
\endhead
\texttt{monotone} & \(\forall t: F^{(t+1)} \leq F^{(t)}\) & FE
decreasing at every step \\
\texttt{total\_reduction} & \(F^{(0)} - F^{(T)}\) & Total free energy
minimised (absolute) \\
\texttt{relative\_reduction\_pct} &
\(100\cdot (F^{(0)} - F^{(T)}) / \lvert F^{(0)}\rvert\) & Fraction of
initial FE eliminated, \% \\
\texttt{n\_iterations} & \(T + 1\) & Number of iterations in the
trajectory \\
\texttt{converged} &
\(\lvert F^{(T)} - F^{(T-1)}\rvert < 0.01\cdot(F_{\max} - F_{\min})\) &
Reached stable minimum (within 1 \% of range) \\
\texttt{fe\_range} & \((F_{\min},\, F_{\max})\) & Absolute FE bounds
across the trajectory \\
\texttt{mean\_step\_size} &
\(\tfrac{1}{T}\sum_{t} \lvert F^{(t+1)} - F^{(t)}\rvert\) & Average
per-iteration FE change \\
\texttt{final\_delta} & \(\lvert F^{(T)} - F^{(T-1)}\rvert\) & Final
step size at convergence or timeout \\
\bottomrule\noalign{}
\end{longtable}

\textbf{Distributional KL divergence} (\texttt{distributional\_kl()})
computes the KL divergence between two discrete return distributions:

\begin{equation}
D_{\mathrm{KL}}(P \| Q) = \sum_{i} P(z_i) \log \frac{P(z_i)}{Q(z_i) + \epsilon}
\label{eq:eq-7c-dkl}
\end{equation}

with \(\epsilon = 10^{-10}\) for numerical stability. Verified
properties: \(D_{\mathrm{KL}}(P\|Q) \geq 0\) (Gibbs' inequality),
\(D_{\mathrm{KL}}(P\|P) = 0\), asymmetry
\(D_{\mathrm{KL}}(P\|Q) \neq D_{\mathrm{KL}}(Q\|P)\) in general.

\textbf{Quantile coverage} (\texttt{quantile\_coverage()}) measures
calibration error---the mean absolute deviation between nominal quantile
levels and empirical coverage frequencies:

\begin{equation}
\mathrm{CE} = \frac{1}{N} \sum_{i=1}^{N} \left| \tau_i - \hat{F}(z_{\tau_i}) \right|
\label{eq:eq-7c-ce}
\end{equation}

A perfectly calibrated distributional model achieves
\(\mathrm{CE} = 0\); the DAIF implementation in
\texttt{src/daif/metrics.py} (tested in \texttt{test\_daif\_metrics.py})
yields \(\mathrm{CE} < 0.05\) on evaluated cases (\textbf{implemented
and tested}).

\textbf{Return distribution entropy}
(\texttt{return\_distribution\_entropy()}) quantifies uncertainty in the
distributional belief:

\begin{equation}
H[Z] = -\sum_{i=1}^{N_{\text{bins}}} p_i \log p_i
\label{eq:eq-7c-entropy}
\end{equation}

where the \(p_i\) are obtained by discretising the
quantile-parameterised return onto \(N_{\text{bins}}\) equal-width bins
(default \(N_{\text{bins}}=50\) in \texttt{src/daif/metrics.py}) with
additive \(\epsilon\)-smoothing
\(p_i \leftarrow (c_i + \epsilon) / (\sum_j c_j + N_{\text{bins}}\,\epsilon)\)
to keep the estimator finite when bins are empty. The estimator is
consistent at the usual \(O(1/N_{\text{bins}})\) discretisation rate as
\(N_{\text{bins}} \to \infty\), and satisfies
\(0 \le H[Z] \le \log N_{\text{bins}}\) by direct enumeration. This
links to the belief trajectory in \autoref{fig:daif-belief-trajectory}:
entropy decreases monotonically as the distributional belief sharpens,
providing a scalar summary of parse certainty.

\begin{block}{Two Supporting Utilities Exposed by \texttt{src/daif/}}
\protect\phantomsection\label{sec:daif-support-utils}
Two further public symbols in the DAIF subpackage support the machinery
above without requiring a separate equation:

\begin{itemize}
\tightlist
\item
  \textbf{\texttt{distributional\_epistemic\_value()}}
  (\texttt{src/daif/policy.py}) measures the information-theoretic value
  of resolving uncertainty in the return distribution via the
  differential-entropy form
  \(\mathrm{EV}_{\mathrm{dist}} = \tfrac{1}{2}\log\bigl(\mathrm{Var}[Z]/\sigma^2_{\mathrm{ref}}\bigr)\).
  In the risk-sensitive regime (\(\beta_{\mathrm{risk}}>0\) in
  \autoref{eq:eq-7c-g}) this function decomposes the risk term
  \(\mathrm{Var}_Z[R(\pi)]\) into a positive ``exploration premium''
  when the current return distribution is more dispersed than the
  reference.
\item
  \textbf{\texttt{categorical\_return\_distribution()}}
  (\texttt{src/daif/core.py}) realises the C51 projection operator
  \(\Phi\colon \mathcal{P}(\mathbb{R}) \to \mathcal{P}(\{z_1,\dots,z_{N_{\text{atoms}}}\})\)
  that maps a quantile-parameterised return onto the fixed atomic
  support of Eq. \ref{eq:eq-7c-c51}. The operator is
  \emph{non-expansive} in \(W_1\): by construction \(\Phi\)
  redistributes each quantile mass across at most two adjacent atoms
  using barycentric weights summing to 1, so for any two return
  distributions \(P, Q\) one has \(W_1(\Phi P, \Phi Q) \le W_1(P, Q)\)
  --- hence \(\Phi\) is 1-Lipschitz (and in fact strictly contractive
  whenever \(P\) and \(Q\) place mass strictly between atoms). This is
  the reason C51/DAIF interoperability preserves contraction bounds from
  \autoref{sec:daif-pushforward}.
\end{itemize}
\end{block}

\begin{block}{Dimensional Analysis}
\protect\phantomsection\label{dimensional-analysis}
The quantities appearing in this section carry the following units:

{\def\LTcaptype{none} % do not increment counter
\begin{longtable}[]{@{}
  >{\raggedright\arraybackslash}p{(\linewidth - 4\tabcolsep) * \real{0.3077}}
  >{\centering\arraybackslash}p{(\linewidth - 4\tabcolsep) * \real{0.3846}}
  >{\raggedright\arraybackslash}p{(\linewidth - 4\tabcolsep) * \real{0.3077}}@{}}
\toprule\noalign{}
\begin{minipage}[b]{\linewidth}\raggedright
Quantity
\end{minipage} & \begin{minipage}[b]{\linewidth}\centering
Symbol
\end{minipage} & \begin{minipage}[b]{\linewidth}\raggedright
Units
\end{minipage} \\
\midrule\noalign{}
\endhead
Variational free energy & \(F\) & nats \\
Return (discounted cumulative reward) & \(Z(s,a)\) & return units
(e.g.~dimensionless log-probability) \\
Wasserstein distance on returns & \(W_p(Z_a, Z_b)\) & return units \\
\(\mathrm{DPE}_{\text{semantic}}\) & --- & return units (=
\(\lvert\Delta\mathbb{E}[Z]\rvert\)) \\
\(\mathrm{DPE}_{\text{structural}}\) & --- & return units (= \(W_1\)) \\
Enriched weight (precision) & \(w_c\) & dimensionless, \(\in[0,1]\) \\
Precision-update magnitude & \(\Delta\Lambda\) & dimensionless (weight
difference) \\
Severity gating & \(S_{\text{violation}}\) & dimensionless,
\(\in\{0,0.5,1\}\) \\
N400 amplitude (\autoref{eq:eq-7c-n400}) & --- & return units
(\emph{not} μV until calibrated) \\
P600 amplitude (\autoref{eq:eq-7c-p600}) & --- & return units ×
dimensionless \(s\) (\emph{not} μV until calibrated) \\
\bottomrule\noalign{}
\end{longtable}
}

Since \(Z\) in our implementation is built from log-likelihood proxies
(\texttt{safe\_log\_lik} in \texttt{src/daif/inference.py}), the return
is effectively in \emph{nats}, and every quantity derived from \(Z\) is
therefore dimensionally a free-energy-like scalar. A future
ERP-calibration step would convert ``nats'' to ``μV'' via a per-subject
scaling constant, as noted in the \autoref{fig:daif-erp-predictions}
caption.
\end{block}
\end{frame}

\begin{frame}[fragile,allowframebreaks]{Limitations and Neurobiological
Scope}
\protect\phantomsection\label{sec:daif-limitations}
Four limitations of the current DAIF implementation are recorded
explicitly for future work:

\begin{enumerate}
\item
  \textbf{Mean-field approximation (cost/accuracy).}
  \texttt{push\_forward\_return()} maintains one belief-weighted return
  distribution rather than per-state distributions \(Z(s)\), reducing
  memory from \(\mathcal{O}(n\cdot N_{\text{atoms}})\) to
  \(\mathcal{O}(n)\). By the mean-field bound proved later in this
  section (the dominated-convergence argument under ``B5'' in the
  contraction analysis) the approximation error is at most
  \(\gamma\cdot R_{\max}\cdot H[q]\) in \(W_1\), where \(R_{\max}\) is
  the sup-norm of the reward vector. For sharp posteriors
  (\(H[q]\lesssim 0.1\) nats at \(\gamma=0.99\) and unit-scale rewards)
  the error is below \(0.1\) return-unit --- well below the
  within-subject noise floor on ERP measurements; it degrades linearly
  as the posterior diffuses.
\item
  \textbf{Enriched-categorical unification is a conjecture.} The three
  ``distributional'' tracks --- distributional semantics
  (\autoref{sec:categorical-semantics}), distributional RL (this
  section), and active-inference posteriors --- are implemented and
  empirically correspond via \([0,1]\)-enriched hom-values, but a
  \emph{categorical} proof that they share a common enriched base
  remains open. The conjecture is stated in the opening of this section
  and is not used as a load-bearing claim elsewhere in the paper.
\item
  \textbf{Empirical validation is narrow.} Our case-assignment
  demonstrations use a single German transitive sentence (\emph{``Der
  Hund jagt die Katze schnell''}). Cross-linguistic and cross-register
  validation is left to future work; the hooks in
  \texttt{make\_daif\_belief\_trajectory\_data()} make adding new
  sentences a single-function change. A Russian or Serbian/BCS sentence
  --- say \emph{Sobaka kusaet čeloveka} ``the dog bites the man'' ---
  would be a particularly clean DAIF stress-test, since the unambiguous
  case suffixes (\emph{-a} NOM.SG.F vs \emph{-a} ACC.SG.M.ANIM after
  stem hardening) deliver an information-theoretically sharper drop in
  \(H[q]\) at the case-marked noun than the German example, where
  word-final case markers compete with positional disambiguation and
  gender / number ambiguity in the determiner system.
\item
  \textbf{Phase--amplitude coupling (PAC) latency gap.} DAIF predicts
  ERP \emph{amplitudes} from distributional prediction errors but does
  not predict component \emph{latencies}. The ROSE model
  {[}@murphy2023rose{]} argues that the structural-discrepancy
  (P600-analogous slow-phase) signal must establish a geometric
  ``mesoscopic protectorate'' before semantic surprise (N400-analogous
  rapid gamma binding) can be fully constrained. A principled treatment
  of this cross-frequency-coupling delay would require an explicit
  timing parameter at the CEREBRUM layer (\autoref{sec:cerebrum}) ---
  the present implementation keeps N400/P600 latencies as fixed Gaussian
  peaks at 380 ms and 600 ms respectively (see
  \texttt{DEFAULT\_N400\_PEAK\_MS} and \texttt{DEFAULT\_P600\_PEAK\_MS}
  in \texttt{src/daif/prediction.py}).
\end{enumerate}
\end{frame}

\begin{frame}[fragile,allowframebreaks]{CEREBRUM: Eight Cases as
Functional Specializations}
\protect\phantomsection\label{sec:cerebrum}
The preceding six contributions---push-forward returns, quantile TD for
case precision, VMP message passing, epistemic policy selection, ERP
convergence profiles, and Bethe free-energy decomposition---define a
distributional active inference layer for sentence processing. CEREBRUM
translates this layer into a complete computational architecture by
assigning each of the eight traditional cases a functional role within
the inference engine.

\begin{block}{Architecture and Design Principles}
\protect\phantomsection\label{architecture-and-design-principles}
The \textbf{CEREBRUM} framework
{[}@friedman2024cerebrum{]}---Case-Enabled Reasoning Engine with
Bayesian Representations for Unified Modeling---provides a computational
architecture that implements the categorical case framework within an
active inference engine. CEREBRUM instantiates the view of Vasil et al.
{[}-@vasil2020world{]} that human communication is itself active
inference: a process of jointly constructing and refining generative
models of shared relational structure.

CEREBRUM's key design principles (\autoref{tbl:cerebrum-principles}):

\begin{longtable}[]{@{}
  >{\raggedright\arraybackslash}p{(\linewidth - 2\tabcolsep) * \real{0.5000}}
  >{\raggedright\arraybackslash}p{(\linewidth - 2\tabcolsep) * \real{0.5000}}@{}}
\caption{CEREBRUM design principles: conceptual commitments and their
implementation in the reasoning
engine.}\label{tbl:cerebrum-principles}\tabularnewline
\toprule\noalign{}
\begin{minipage}[b]{\linewidth}\raggedright
Principle
\end{minipage} & \begin{minipage}[b]{\linewidth}\raggedright
Implementation
\end{minipage} \\
\midrule\noalign{}
\endfirsthead
\toprule\noalign{}
\begin{minipage}[b]{\linewidth}\raggedright
Principle
\end{minipage} & \begin{minipage}[b]{\linewidth}\raggedright
Implementation
\end{minipage} \\
\midrule\noalign{}
\endhead
\textbf{Cases as functional roles} & Model components carry case
markings that determine their computational role in the inference
cycle \\
\textbf{Morphisms as message passing} & Grammatical relations are
implemented as message-passing channels between components \\
\textbf{Enriched weights as precision} & The \([0,1]\) weights on
morphisms correspond to precision parameters in the variational
inference scheme \\
\textbf{Alignment as model selection} & Different alignment types
correspond to different generative model architectures, selected by
Bayesian model comparison \\
\textbf{Diagrams as generative models} & Commutative diagrams serve as
the structural specification of the generative model \\
\textbf{DAIF as distributional layer} & The \texttt{src/daif/}
subpackage provides the distributional RL layer: full return
distributions replace point estimates throughout the generative cycle \\
\bottomrule\noalign{}
\end{longtable}
\end{block}

\begin{block}{Case Roles as Functional Specializations in CEREBRUM}
\protect\phantomsection\label{case-roles-as-functional-specializations-in-cerebrum}
CEREBRUM deploys the eight traditional cases as functional
specializations, each with a DAIF-level extension
(\autoref{tbl:cerebrum-daif}):

\begin{table}[htbp]
\centering
\begin{tabular}{@{}l p{0.29\textwidth} p{0.29\textwidth} p{0.29\textwidth}@{}}
\toprule
\textbf{Case} & \textbf{CEREBRUM Function} & \textbf{Active Inference Role} & \textbf{DAIF Extension} \\
\midrule
NOM & Primary driver / agent & Source of action policies & Softmax policy over $G{(\pi)}$; highest epistemic value \\
ACC & Primary target / patient & Object of predictions & Predicted distribution $Z_{\text{acc}}$; error-driven update \\
GEN & Source / possessor & Provider of priors & Prior return distribution $p(Z)$ \\
DAT & Recipient / goal & Target of information transfer & EIG maximised toward DAT state \\
INS & Instrument / means & Tool for state transformation & IQN risk distortion (neutral mode) \\
LOC & Context / environment & Markov blanket boundary & Bethe FE boundary conditions \\
ABL & Origin / cause & Source of causal influence & Push-forward source measure $\mathbb{P}_{x_0,a_0}$ \\
VOC & Addressee & Pragmatic pointer & Lowest epistemic weight; largest P600 on violation \\
\bottomrule
\end{tabular}
\caption{CEREBRUM case roles as functional specializations with DAIF distributional extensions.}
\label{tbl:cerebrum-daif}
\end{table}
\end{block}
\end{frame}

\end{document}
